% GENERATED from QPf2-draw-attach.md by md2tex.py. Do not hand-edit: sync is
% one-way and the next run overwrites this file.

\section{Every drawing has one file, and one owner}
\textbf{The files}: one drawing per owner, each in its owner's folder.

\begin{verbatim}
 page drawing     PAGEDIR/draw/PAGEID.excalidraw    the page's own scene
 group drawing    GROUPDIR/draw/group.excalidraw    layout + relations over the pages
 board map        board.excalidraw                  the Index page's overview scene
 image bytes      draw/assets/OWNER/                pasted images, kept out of the scene
\end{verbatim}
📌 This part says which file holds which drawing, so you always know what to open. The two page names in that list are different on purpose. The folder is named for the page id plus its slug. The scene file is named for the page id alone. This page sits in \texttt{QPf2-draw-attach/} and its scene is \texttt{draw/QPf2.excalidraw}. Each page drawing is a plain Excalidraw file, and it is still useful opened on its own. A pasted image is saved under the owner's \texttt{draw/assets/}, and the scene keeps a relative path to it. A page with no drawing yet gets an empty scene the first time its Draw split opens. So there is never a missing file to trip over. The group drawing holds only the group's own layout and a list naming each page's scene. It never copies what those scenes contain. \texttt{board.excalidraw} stays at the board root as the Index page's overview map. Everything below it follows the new layout: one scene per page, one per group.


\section{You draw beside the page, never inside it}
\textbf{The Draw split}: the drawing opens next to the page, never inside it.

\begin{verbatim}
 open        Plugin   Draw, or the Draw tab in the right pane
 edit        draw on the canvas; saves go to this page's own scene file
 generate    the Draw-it button: Claude draws the page's ## Diagram for you
 follow      navigate to another page and the canvas follows it
\end{verbatim}
📌 This part says how to open a drawing, how to edit it, and how to ask Claude to draw it for you. The split names the file it will save to before you draw. So a stroke can never land in another page's scene. The Draw-it button takes an optional ask. Leave it empty, and it turns the page's ascii figure into a real drawing. A drawing Claude made can be drawn again as often as you like. The button refuses to overwrite a drawing a person made by hand.


\section{The big group picture is built from the page files, not copied from them}
\textbf{One large drawing from many small ones}: what is kept, and what is built for you.

\begin{verbatim}
 kept       each page's scene . the group's own layer . where each one goes
 built      one group view, built again on every open
 owner      every shape keeps its owner, so an edit goes back to the right file
\end{verbatim}
📌 This part says the group picture is put together from the page files each time it opens. Opening a group's drawing loads the group scene, pulls in every member page's scene, and places them where the layout says. Inside the group view you work in one mode at a time. You can arrange the pages, edit the group's own layer, or step into one page's scene. Only that owner's file receives the save. Save a page, and the group view shows the change the next time it opens. Nothing is ever merged or pulled in by hand.


\section{Two people cannot wipe out each other's work}
\textbf{The safe save}: what the server checks before it writes.

\begin{verbatim}
 opened     the revision the editor loaded
 save       owner . that revision . the changed elements
[x] accepted   the file has not moved since
 conflict   someone else saved first . reload and compare
\end{verbatim}
📌 This part says every save is checked first, so nobody quietly loses work. Every save carries the version it started from. So two people on the same drawing cannot quietly overwrite each other. The chat pane gets the same owner address as the canvas. So asking chat to change a drawing writes to that drawing's own file.

